% Registered abstract of HotCRP paper 579, verbatim (only the math delimiters of $\ell_2$ are LaTeX).
% SaTML forbids substantial changes after registration; the post-hoc two-sided result is stated in the
% Introduction (Findings) and Section~\ref{sec:grclip}. Re-add a clarifying clause only with written chair approval.
Randomized smoothing can provide certified $\ell_2$ robustness for a smoothed classifier, but a stable prediction need not be a correct one. We study whether post-hoc shared corrections of vision-language prototype geometry improve the certified correctness of frozen zero-shot CLIP classifiers under pixel-space Gaussian smoothing. We first show, through explicit normalized constructions, that a clean global modality-gap summary does not by itself determine the local decision geometry or noisy class probabilities that govern a Cohen certificate. We then derive exact per-draw identities and necessary decision-change conditions for the studied shared prototype operators, while keeping these geometric diagnostics separate from the randomized-smoothing certificate itself. In a preregistered, custody-controlled audit across frozen CLIP backbones and datasets, the registered family of shared/global corrections did not establish an improvement in standard Cohen certified accuracy. In several difficult regimes, raw-clean and smoothed predictions disagree substantially, and certified-wrong predictions concentrate on a small number of classes; we report these observations descriptively rather than as identified causes. We additionally evaluate a separate supervised class-specific positive control to determine whether certified correctness remains recoverable outside the tested shared operator family. We will release registrations, sufficient statistics, and reconstruction code, and argue that evaluations of certified vision-language models should report correctness, stability, abstention, and certified-wrong concentration separately.
